\documentclass[../../main.tex]{subfiles}
\begin{document}

\section{切丛}\label{section:ch03-切丛}

常常有用的是考虑一个流形所有点处所有切向量所成的集合。给定一个带边或无边的光滑流形 $M$，我们定义 $M$ 的切丛，记作 $TM$，为 $M$ 所有点处切空间的不交并：
\begin{align*}
TM=\bigsqcup_{p\in M}T_{p}M.
\end{align*}
我们通常把这个不交并的元素写成有序对 $(p,v)$，其中 $p\in M$、$v\in T_{p}M$（而不是把点 $p$ 放在第二个位置，像不交并元素更常见的写法那样）。切丛自然带有一个投影映射 $\pi:TM\to M$，它把 $T_{p}M$ 中每个向量送到它所切于的点 $p$：$\pi(p,v)=p$。我们将常常犯下那种常见的轻微罪过，把 $T_{p}M$ 与其在典范单射 $v\mapsto(p,v)$ 下的像等同，并将根据我们希望给点 $p$ 多少强调，对 $T_{p}M$ 中的一个切向量任意使用记号 $(p,v)$、$v_{p}$ 和 $v$。

例如，在特殊情形 $M=\mathbb{R}^{n}$，利用命题~\cref{proposition:ch03-导子与几何切向量}，我们看到 $\mathbb{R}^{n}$ 的切丛可以典范地等同于其几何切空间的并集，而后者又正是 $\mathbb{R}^{n}$ 与自身的笛卡儿积：
\begin{align*}
T\mathbb{R}^{n}=\bigsqcup_{a\in\mathbb{R}^{n}}T_{a}\mathbb{R}^{n}\cong\bigsqcup_{a\in\mathbb{R}^{n}}\mathbb{R}^{n}_{a}=\bigsqcup_{a\in\mathbb{R}^{n}}(\{a\}\times\mathbb{R}^{n})=\mathbb{R}^{n}\times\mathbb{R}^{n}.
\end{align*}
这个笛卡儿积的元素 $(a,v)$ 可以看作代表几何切向量 $v_{a}$，或由~\eqref{align:ch03-方向导数} 定义的导子 $D_{v|_{a}}$。然而要注意，一般而言，一个光滑流形的切丛不能以任何自然方式与一个笛卡儿积等同，因为没有把不同点处切空间彼此等同的典范方式。我们下面将对此多说几句。

若 $M$ 是一个光滑流形，切丛 $TM$ 可以简单地看作向量空间的一个不交并；但它远不止于此。下一个命题表明 $TM$ 本身可以被看作一个光滑流形。

\begin{proposition}\label{proposition:ch03-切丛流形}
对任意光滑 $n$ 流形 $M$，切丛 $TM$ 有一个自然拓扑和光滑结构，使之成为一个 $2n$ 维光滑流形。关于这个结构，投影 $\pi:TM\to M$ 光滑。
\end{proposition}

\begin{proof}
我们从定义将成为我们光滑坐标图的映射开始。给定 $M$ 的任意光滑坐标图 $(U,\varphi)$，注意 $\pi^{-1}(U)\subset TM$ 是 $M$ 在 $U$ 所有点处的所有切向量所成集合。令 $(x^{1},\dots,x^{n})$ 记 $\varphi$ 的坐标函数，定义一个映射 $\widetilde{\varphi}:\pi^{-1}(U)\to\mathbb{R}^{2n}$ 为
\begin{align}\label{align:ch03-切丛自然坐标}
\widetilde{\varphi}\left(v^{i}\left.\frac{\partial}{\partial x^{i}}\right|_{p}\right)=\bigl(x^{1}(p),\dots,x^{n}(p),v^{1},\dots,v^{n}\bigr).
\end{align}
（见图~\cref{figure:ch03-3.8}。）它的像集是 $\varphi(U)\times\mathbb{R}^{n}$，这是 $\mathbb{R}^{2n}$ 的一个开子集。它是到其像的双射，因为其逆可以显式写成
\begin{align*}
\widetilde{\varphi}^{-1}(x^{1},\dots,x^{n},v^{1},\dots,v^{n})=v^{i}\left.\frac{\partial}{\partial x^{i}}\right|_{\varphi^{-1}(x)}.
\end{align*}

\begin{figure}[ht]
\centering
\includegraphics[width=0.8\textwidth]{fig-03-8.png}
\caption{切丛的坐标}\label{figure:ch03-3.8}
\end{figure}

现在设给定 $M$ 的两个光滑坐标图 $(U,\varphi)$ 和 $(V,\psi)$，令 $(\pi^{-1}(U),\widetilde{\varphi})$、$(\pi^{-1}(V),\widetilde{\psi})$ 为 $TM$ 上相应的坐标图。集合
\begin{align*}
\widetilde{\varphi}\bigl(\pi^{-1}(U)\cap\pi^{-1}(V)\bigr)=\varphi(U\cap V)\times\mathbb{R}^{n}
\end{align*}
和
\begin{align*}
\widetilde{\psi}\bigl(\pi^{-1}(U)\cap\pi^{-1}(V)\bigr)=\psi(U\cap V)\times\mathbb{R}^{n}
\end{align*}
在 $\mathbb{R}^{2n}$ 中是开的，且过渡映射 $\widetilde{\psi}\circ\widetilde{\varphi}^{-1}:\varphi(U\cap V)\times\mathbb{R}^{n}\to\psi(U\cap V)\times\mathbb{R}^{n}$ 可以用~\eqref{align:ch03-分量变换} 显式写成
\begin{align*}
\widetilde{\psi}\circ\widetilde{\varphi}^{-1}(x^{1},\dots,x^{n},v^{1},\dots,v^{n})
=\left(\widetilde{x}^{1}(x),\dots,\widetilde{x}^{n}(x),\frac{\partial\widetilde{x}^{1}}{\partial x^{j}}(x)v^{j},\dots,\frac{\partial\widetilde{x}^{n}}{\partial x^{j}}(x)v^{j}\right).
\end{align*}
这显然光滑。

选取 $M$ 的一个由光滑坐标域组成的可数覆盖 $\{U_{i}\}$，我们得到 $TM$ 的一个由坐标域 $\{\pi^{-1}(U_{i})\}$ 组成的可数覆盖，它满足光滑流形坐标图引理（引理~1.35）的条件~(i)--(iv)。为验证 Hausdorff 条件~(v)，只需注意同一根 $\pi$ 的纤维中任意两点落在一个坐标图中，而若 $(p,v)$ 和 $(q,w)$ 落在不同纤维中，则存在 $M$ 的不相交光滑坐标域 $U,V$ 使 $p\in U$、$q\in V$，于是 $\pi^{-1}(U)$ 和 $\pi^{-1}(V)$ 是分别包含 $(p,v)$ 和 $(q,w)$ 的不相交坐标邻域。

为看出 $\pi$ 光滑，注意关于 $M$ 的坐标图 $(U,\varphi)$ 和 $TM$ 的坐标图 $(\pi^{-1}(U),\widetilde{\varphi})$，它的坐标表示是 $\pi(x,v)=x$。
\end{proof}

由~\eqref{align:ch03-切丛自然坐标} 给出的坐标 $(x^{i},v^{i})$ 称为 $TM$ 上的自然坐标。

\begin{exercise}\label{exercise:ch03-3.19}
设 $M$ 是一个光滑带边流形。证明 $TM$ 有一个自然拓扑和光滑结构，使之成为一个光滑带边流形，使得若 $(U,(x^{i}))$ 是 $M$ 的任意光滑边界坐标图，则重排 $TM$ 的自然坐标图 $(\pi^{-1}(U),(x^{i},v^{i}))$ 中的坐标，便得到一个边界坐标图 $(\pi^{-1}(U),(v^{i},x^{i}))$。
\end{exercise}

\begin{proposition}\label{proposition:ch03-单图覆盖切丛}
若 $M$ 是一个光滑 $n$ 维带边或无边流形，且 $M$ 可被单个光滑坐标图覆盖，则 $TM$ 微分同胚于 $M\times\mathbb{R}^{n}$。
\end{proposition}

\begin{proof}
若 $(U,\varphi)$ 是 $M$ 的一个整体光滑坐标图，则特别地 $\varphi$ 是从 $U=M$ 到开子集 $\widetilde{U}\subset\mathbb{R}^{n}$ 或 $\mathbb{H}^{n}$ 的微分同胚。前一命题的证明表明自然坐标图 $\widetilde{\varphi}$ 是从 $TM$ 到 $\widetilde{U}\times\mathbb{R}^{n}$ 的双射，且 $TM$ 上的光滑结构本质上是通过宣布 $\widetilde{\varphi}$ 为微分同胚来定义的。
\end{proof}

虽然乘积 $U\times\mathbb{R}^{n}$ 的图景是局部可视化切丛上光滑结构的有用方式（如图~\cref{figure:ch03-3.8}），但不要被误导而以为每个切丛都整体微分同胚（甚至同胚）于该流形与 $\mathbb{R}^{n}$ 的乘积。对大多数光滑流形而言并非如此。我们将在第~8、10、16 章重新讨论这个问题。

通过把 $F$ 在 $M$ 所有点处的微分合在一起，我们得到切丛之间一个整体定义的映射，称为整体微分或整体切映射，记作 $\mathrm{d}F:TM\to TN$。这正是那个在每个切空间 $T_{p}M\subset TM$ 上的限制为 $\mathrm{d}F_{p}$ 的映射。当我们把 $F$ 的微分应用于一个具体向量 $v\in T_{p}M$ 时，根据我们希望给点 $p$ 多少强调，我们可以写 $\mathrm{d}F_{p}(v)$ 或 $\mathrm{d}F(v)$。前者信息量更大，而后者更简洁。

我们在 $TM$ 上定义的光滑结构的一个重要特征是，它使一个光滑映射的微分成为切丛之间的一个光滑映射。

\begin{proposition}\label{proposition:ch03-整体微分光滑}
若 $F:M\to N$ 是一个光滑映射，则其整体微分 $\mathrm{d}F:TM\to TN$ 是一个光滑映射。
\end{proposition}

\begin{proof}
由 $\mathrm{d}F_{p}$ 在坐标中的局部表达式~\eqref{align:ch03-欧氏微分基}，可得 $\mathrm{d}F$ 关于 $TM$ 和 $TN$ 的自然坐标有如下坐标表示：
\begin{align*}
\mathrm{d}F(x^{1},\dots,x^{n},v^{1},\dots,v^{n})
=\left(F^{1}(x),\dots,F^{n}(x),\frac{\partial F^{1}}{\partial x^{i}}(x)v^{i},\dots,\frac{\partial F^{n}}{\partial x^{i}}(x)v^{i}\right).
\end{align*}
这是光滑的，因为 $F$ 光滑。
\end{proof}

整体微分的下列性质直接由命题~\cref{proposition:ch03-微分性质} 得出。

\begin{corollary}[整体微分的性质]\label{corollary:ch03-整体微分性质}
设 $F:M\to N$ 和 $G:N\to P$ 是光滑映射。
\begin{enumerate}[(a)]
\item $\mathrm{d}(G\circ F)=\mathrm{d}G\circ\mathrm{d}F$。
\item $\mathrm{d}(\mathrm{Id}_{M})=\mathrm{Id}_{TM}$。
\item 若 $F$ 是微分同胚，则 $\mathrm{d}F:TM\to TN$ 也是微分同胚，且 $(\mathrm{d}F)^{-1}=\mathrm{d}(F^{-1})$。
\end{enumerate}
\end{corollary}

由于这个推论的~(c) 部分，当 $F$ 是微分同胚时，我们可以无歧义地使用记号 $\mathrm{d}F^{-1}$ 来表示 $(\mathrm{d}F)^{-1}$ 或 $\mathrm{d}(F^{-1})$。

\end{document}